% GENERATED FILE. Edit canonical lessons, then run npm run generate:books.
% canonical-source-hash: fnv1a64:5cb4ac0fe8f36b90
% canonical-lessons: JA-C188-yuubinkyoku, JA-C188-kitte, JA-C188-hagaki, JA-C188-okuru, JA-C188-kakitome

\chapter{At the Post Office}
\label{ch:ja-at-the-post-office}
\hlchaptermodality{188}

\begin{chapteropening}
\textbf{By the end of this chapter:} \emph{I can buy stamps and send post at the post office.}

Complete the last lesson of chapter 188: I can buy stamps and send post at the post office.
\end{chapteropening}

\section[\texorpdfstring{\ja{ゆうびんきょく} (yūbinkyoku)}{ゆうびんきょく (yūbinkyoku)}]{\ja{ゆうびんきょく} (yūbinkyoku) — post office}

\label{lesson:JA-C188-yuubinkyoku}



\begin{quote}
Before the new one: say the Japanese for to borrow, then the Japanese for money exchange.
\end{quote}

\subsection*{You'll want to know: \ja{ゆうびんきょく}}
\textbf{\ja{ゆうびんきょく}} — \emph{yūbinkyoku} — \textquotedblleft{}post office\textquotedblright{}.

\begin{quote}
\textbf{\ja{ゆうびんきょくは} \ja{えきの} \ja{ちかくです。}} — \emph{yūbinkyoku wa eki no chikaku desu} — \textquotedblleft{}The post office is near the station.\textquotedblright{}
\end{quote}

\begin{cousinweb}
\textbf{\ja{ゆうびんきょく}} (\emph{yūbinkyoku}) is Sino-Japanese (kango): three Chinese characters read with their on'yomi, readings borrowed from Chinese. They are \emph{yū}, \textquotedblleft{}post station\textquotedblright{} (Middle Chinese \emph{hjuw}), \emph{bin}, \textquotedblleft{}convenience\textquotedblright{} (\emph{bjienH}) and \emph{kyoku}, \textquotedblleft{}office, bureau\textquotedblright{} (\emph{gjowk}) (KANJIDIC2; Middle Chinese in Baxter's transcription).
\end{cousinweb}

\begin{grammarlens}[title={What you've built}]
One more word for everyday life at A2.
\end{grammarlens}

\subsection*{Your turn}
\begin{itemize}
\raggedright
  \item \emph{Say it:} \emph{yūbinkyoku}
  \item \emph{Say it:} \emph{yūbinkyoku}, once more
  \item \emph{Say it:} \emph{yūbinkyoku wa eki no chikaku desu} — The post office is near the station.
  \item \emph{Recall:} say the Japanese for to borrow, then the Japanese for money exchange, then say \emph{yūbinkyoku} again
\end{itemize}

\begin{tcolorbox}[breakable,colback=teal!4,colframe=teal!35!black,title={Before you move on}]
What does \ja{ゆうびんきょく} mean? (\textquotedblleft{}Post office\textquotedblright{}.) Say it once more.
\end{tcolorbox}
\section[\texorpdfstring{\ja{きって} (kitte)}{きって (kitte)}]{\ja{きって} (kitte) — stamp; postage stamp}

\label{lesson:JA-C188-kitte}



\begin{quote}
Before the new one: say the Japanese for money exchange, then the Japanese for post office.
\end{quote}

\subsection*{You'll want to know: \ja{きって}}
\textbf{\ja{きって}} — \emph{kitte} — \textquotedblleft{}stamp; postage stamp\textquotedblright{}.

\begin{quote}
\textbf{\ja{ゆうびんきょくで} \ja{きってを} \ja{かいます。}} — \emph{yūbinkyoku de kitte o kaimasu} — \textquotedblleft{}I buy stamps at the post office.\textquotedblright{}
\end{quote}

\begin{cousinweb}
\textbf{\ja{きって}} (\emph{kitte}) is native Japanese, built from \emph{kiri}, the stem of \emph{kiru}, \textquotedblleft{}to cut\textquotedblright{}, and \emph{te}, \textquotedblleft{}hand\textquotedblright{}, both inherited rather than borrowed (JMdict, KANJIDIC2).
\end{cousinweb}

\begin{grammarlens}[title={What you've built}]
One more word for everyday life at A2.
\end{grammarlens}

\subsection*{Your turn}
\begin{itemize}
\raggedright
  \item \emph{Say it:} \emph{kitte}
  \item \emph{Say it:} \emph{kitte}, once more
  \item \emph{Say it:} \emph{yūbinkyoku de kitte o kaimasu} — I buy stamps at the post office.
  \item \emph{Recall:} say the Japanese for money exchange, then the Japanese for post office, then say \emph{kitte} again
\end{itemize}

\begin{tcolorbox}[breakable,colback=teal!4,colframe=teal!35!black,title={Before you move on}]
What does \ja{きって} mean? (\textquotedblleft{}Stamp; postage stamp\textquotedblright{}.) Say it once more.
\end{tcolorbox}
\section[\texorpdfstring{\ja{はがき} (hagaki)}{はがき (hagaki)}]{\ja{はがき} (hagaki) — postcard}

\label{lesson:JA-C188-hagaki}



\begin{quote}
Before the new one: say the Japanese for post office, then the Japanese for stamp.
\end{quote}

\subsection*{You'll want to know: \ja{はがき}}
\textbf{\ja{はがき}} — \emph{hagaki} — \textquotedblleft{}postcard\textquotedblright{}.

\begin{quote}
\textbf{\ja{ともだちに} \ja{はがきを} \ja{かきます。}} — \emph{tomodachi ni hagaki o kakimasu} — \textquotedblleft{}I write a postcard to a friend.\textquotedblright{}
\end{quote}

\begin{cousinweb}
\textbf{\ja{はがき}} (\emph{hagaki}) is native Japanese, built from \emph{ha} and \emph{kaki}, \textquotedblleft{}writing\textquotedblright{}, from \emph{kaku}, \textquotedblleft{}to write\textquotedblright{}; JMdict spells the first part as either \textquotedblleft{}leaf\textquotedblright{} or \textquotedblleft{}edge\textquotedblright{}, so which came first is not settled (JMdict, KANJIDIC2).
\end{cousinweb}

\begin{grammarlens}[title={What you've built}]
One more word for everyday life at A2.
\end{grammarlens}

\subsection*{Your turn}
\begin{itemize}
\raggedright
  \item \emph{Say it:} \emph{hagaki}
  \item \emph{Say it:} \emph{hagaki}, once more
  \item \emph{Say it:} \emph{tomodachi ni hagaki o kakimasu} — I write a postcard to a friend.
  \item \emph{Recall:} say the Japanese for post office, then the Japanese for stamp, then say \emph{hagaki} again
\end{itemize}

\begin{tcolorbox}[breakable,colback=teal!4,colframe=teal!35!black,title={Before you move on}]
What does \ja{はがき} mean? (\textquotedblleft{}Postcard\textquotedblright{}.) Say it once more.
\end{tcolorbox}
\section[\texorpdfstring{\ja{おくる} (okuru)}{おくる (okuru)}]{\ja{おくる} (okuru) — to send; to see off}

\label{lesson:JA-C188-okuru}



\begin{quote}
Before the new one: say the Japanese for stamp, then the Japanese for postcard.
\end{quote}

\subsection*{You'll want to know: \ja{おくる}}
\textbf{\ja{おくる}} — \emph{okuru} — \textquotedblleft{}to send; to see off\textquotedblright{}.

\begin{quote}
\textbf{\ja{ははに} \ja{てがみを} \ja{おくります。}} — \emph{haha ni tegami o okurimasu} — \textquotedblleft{}I send a letter to my mother.\textquotedblright{}
\end{quote}

Polite \textbf{\ja{おくります}}, \emph{okurimasu}; a godan verb.

\begin{cousinweb}
\textbf{\ja{おくる}} (\emph{okuru}) is a native Japanese word (wago), inherited rather than borrowed from Chinese; its kanji was attached later as a kun'yomi (JMdict, KANJIDIC2).
\end{cousinweb}

\begin{grammarlens}[title={What you've built}]
One more word for everyday life at A2.
\end{grammarlens}

\subsection*{Your turn}
\begin{itemize}
\raggedright
  \item \emph{Say it:} \emph{okuru}
  \item \emph{Say it:} \emph{okuru}, once more
  \item \emph{Say it:} \emph{haha ni tegami o okurimasu} — I send a letter to my mother.
  \item \emph{Recall:} say the Japanese for stamp, then the Japanese for postcard, then say \emph{okuru} again
\end{itemize}

\begin{tcolorbox}[breakable,colback=teal!4,colframe=teal!35!black,title={Before you move on}]
What does \ja{おくる} mean? (\textquotedblleft{}To send; to see off\textquotedblright{}.) Say it once more.
\end{tcolorbox}
\section[\texorpdfstring{\ja{かきとめ} (kakitome)}{かきとめ (kakitome)}]{\ja{かきとめ} (kakitome) — registered mail; registered post}

\label{lesson:JA-C188-kakitome}



\begin{quote}
Before the new one: say the Japanese for postcard, then the Japanese for to send.
\end{quote}

\subsection*{You'll want to know: \ja{かきとめ}}
\textbf{\ja{かきとめ}} — \emph{kakitome} — \textquotedblleft{}registered mail; registered post\textquotedblright{}.

\begin{quote}
\textbf{\ja{ゆうびんきょくで} \ja{かきとめを} \ja{おくります。}} — \emph{yūbinkyoku de kakitome o okurimasu} — \textquotedblleft{}I send registered mail at the post office.\textquotedblright{}
\end{quote}

\begin{cousinweb}
\textbf{\ja{かきとめ}} (\emph{kakitome}) is native Japanese, built from \emph{kaki}, the stem of \emph{kaku}, \textquotedblleft{}to write\textquotedblright{}, and \emph{tome}, the stem of \emph{tomeru}, \textquotedblleft{}to fasten, to keep in place\textquotedblright{}, both inherited rather than borrowed (JMdict, KANJIDIC2).
\end{cousinweb}

\begin{grammarlens}[title={What you've built}]
One more word for everyday life at A2.
\end{grammarlens}

\subsection*{Your turn}
\begin{itemize}
\raggedright
  \item \emph{Say it:} \emph{kakitome}
  \item \emph{Say it:} \emph{kakitome}, once more
  \item \emph{Say it:} \emph{yūbinkyoku de kakitome o okurimasu} — I send registered mail at the post office.
  \item \emph{Recall:} say the Japanese for postcard, then the Japanese for to send, then say \emph{kakitome} again
\end{itemize}

\begin{tcolorbox}[breakable,colback=teal!4,colframe=teal!35!black,title={Before you move on}]
What does \ja{かきとめ} mean? (\textquotedblleft{}Registered mail; registered post\textquotedblright{}.) Say it once more.
\end{tcolorbox}
